\chapter{Vitamin B12 (Cobalamin)}

\section*{What Is This Test?}
The vitamin B12 test measures the serum concentration of cobalamin, a water-soluble B-complex vitamin essential for DNA synthesis, red blood cell production, and myelin formation. Vitamin B12 is a cobalt-containing corrin ring compound that functions as a coenzyme in two critical reactions: methylcobalamin, which transfers methyl groups in the methionine synthase reaction converting homocysteine to methionine, and adenosylcobalamin, which is required for the conversion of methylmalonyl-CoA to succinyl-CoA. Because folate and B12 metabolism are linked through the methionine cycle, B12 deficiency produces megaloblastic changes identical to folate deficiency. Vitamin B12 is obtained exclusively from animal sources (meat, fish, eggs, dairy), requires gastric acid to be released from food proteins, and is absorbed in the terminal ileum bound to intrinsic factor produced by gastric parietal cells. Deficiency develops slowly (often over years) because total body stores are large (2--5 mg) relative to daily requirements (2.4 mcg). Serum cobalamin is the standard first-line screening test, but because approximately 20--40\% of patients with clinical deficiency have borderline normal serum levels, additional markers (methylmalonic acid and homocysteine) are used to confirm deficiency in ambiguous cases.

\section*{Alternative Names}
Cobalamin; Vitamin B12 Level; Serum B12; Cyanocobalamin; Hydroxocobalamin; Holotranscobalamin (Active B12, holoTC); Total Vitamin B12

\section*{Principle}
Serum vitamin B12 is measured by competitive immunoassay using a chemiluminescent or electrochemiluminescent label (e.g., on Roche, Siemens, Abbott, and Beckman Coulter platforms). In these competitive assays, the patient's B12 competes with a labeled B12 analog for binding to a B12-binding protein (intrinsic factor or a purified B12-binding protein); the signal is inversely proportional to the B12 concentration. Automated assays are standardized against the WHO International Standard for B12, but significant inter-assay variability persists, so results should be interpreted against the reference range of the reporting laboratory. Active B12 (holotranscobalamin, holoTC) is measured by a separate immunoassay that captures B12 bound to transcobalamin, the fraction available for cellular uptake; it becomes abnormal earlier than total B12 and is increasingly used as a first-line marker. Confirmatory markers are measured by alternative methods: methylmalonic acid (MMA) by LC-MS/MS or GC-MS, and homocysteine by immunoassay or HPLC.

\section*{History}
Pernicious anemia was first described by Thomas Addison in 1849 and characterized as a megaloblastic anemia by Biermer in 1872. In 1926, George Minot and William Murphy demonstrated that feeding large quantities of raw liver cured pernicious anemia, earning them (with George Whipple) the 1934 Nobel Prize in Physiology or Medicine. The active factor, vitamin B12, was isolated as cyanocobalamin by Karl Folkers and colleagues in 1948 and by E. Lester Smith independently. The role of intrinsic factor in ileal absorption was established by William Castle in the 1930s. The Schilling test, introduced in 1953, used radiolabeled B12 to distinguish malabsorption due to intrinsic factor deficiency from ileal disease. Radioisotope dilution assays for serum B12 were developed in the 1960s and replaced the microbiological Lactobacillus assay (a growth-based assay first used in 1948). The first automated immunoassay for B12 became available in the 1980s, and the Schilling test was largely abandoned because of the availability of reliable serum markers (B12, MMA, homocysteine) and the disappearance of the radiolabeled source.

\section*{Why This Test Is Ordered}
\begin{itemize}
  \item Evaluation of macrocytic anemia (MCV greater than 100 fL) with possible megaloblastic changes on the peripheral smear (macrocytosis, hypersegmented neutrophils, oval macrocytes)
  \item Investigation of unexplained fatigue, weakness, or cognitive decline in the elderly, in whom B12 deficiency is common and frequently overlooked
  \item Evaluation of peripheral neuropathy, paresthesias, ataxia, or subacute combined degeneration of the spinal cord (the neurological manifestations of B12 deficiency, which can occur without anemia)
  \item Assessment of patients at high risk for deficiency: strict vegetarians and vegans, patients with pernicious anemia, gastric or ileal surgery, Crohn disease, celiac disease, atrophic gastritis, long-term proton pump inhibitor or metformin use, and the elderly (food-cobalamin malabsorption)
  \item Screening and monitoring in patients with a history of bariatric surgery, especially gastric bypass, who require lifelong B12 monitoring
  \item Evaluation of hyperhomocysteinemia as a cardiovascular risk marker; B12, folate, and B6 deficiency are reversible causes
  \item Investigation of infertility, recurrent pregnancy loss, or neural tube defect risk in pregnancy
  \item Monitoring of B12 supplementation in deficient patients to confirm response
  \item Workup of methylmalonic acidemia and other inborn errors of metabolism in infants
\end{itemize}

\section*{Primary vs Secondary Test}
Serum B12 is a primary screening test for B12 deficiency. Because serum B12 has limited sensitivity and specificity at borderline values, the active form (holotranscobalamin) and the functional markers methylmalonic acid and homocysteine serve as secondary confirmatory tests when the serum level is borderline or when clinical suspicion is high despite a normal serum B12.

\section*{How This Test Is Performed}
A blood sample is collected by standard venipuncture into a serum separator tube (gold-top) or plain red-top tube. Fasting is not required. The sample is centrifuged and serum is analyzed on an automated immunoassay platform. For active B12 (holotranscobalamin), a separate immunoassay is used. When confirmatory testing is needed, additional tubes are sent for methylmalonic acid (LC-MS/MS) and homocysteine. Results are typically available within 1--2 days for serum B12; MMA and homocysteine are frequently send-out tests with a 3--7 day turnaround.

\section*{What Preparation Is Needed}
No fasting or special preparation is required. Patients should fast only if B12 is measured as part of a combined panel requiring fasting (e.g., lipids). Because recent oral B12 supplementation or a recent B12 injection rapidly raises serum B12 and can mask true deficiency, patients should disclose supplementation history; ideally the baseline sample is drawn before starting treatment. Intramuscular B12 injections raise serum levels for weeks, so samples should be drawn as close to the next scheduled dose as possible if monitoring is required.

\section*{Contraindications and Precautions}
There are no absolute contraindications to standard venipuncture. Important interpretive cautions: (1) Serum B12 has limited sensitivity --- up to 40\% of patients with clinical B12 deficiency confirmed by elevated MMA have serum B12 in the low-normal range (200--300 pg/mL); in such cases, measure MMA and homocysteine. (2) False-normal serum B12 occurs in myeloproliferative disorders (elevated transcobalamin), liver disease, and autoimmune disease because binding proteins are elevated. (3) False-low serum B12 occurs in pregnancy, oral contraceptive use, folate deficiency, and multiple myeloma due to reduced binding proteins. (4) Severe folate deficiency can cause a functional B12 deficiency picture with elevated MMA. (5) Biotin supplementation (greater than 5 mg/day) interferes with biotin-streptavidin-based immunoassays, producing falsely low or high B12 results depending on the assay --- biotin should be withheld 48 hours before testing. (6) Recent blood transfusion, high-dose vitamin C, and certain medications can interfere with some assays. (7) Gastrectomy and terminal ileal resection permanently impair B12 absorption and mandate lifelong parenteral supplementation regardless of the initial result.

\section*{Risks}
Minimal risk. Slight pain or bruising at the venipuncture site. Rarely: hematoma, infection, or vasovagal syncope.

\section*{What Happens After the Test}
Results are available within 1--7 days depending on whether confirmatory markers are required. A clearly low serum B12 with compatible clinical findings confirms deficiency and is treated with B12 supplementation: parenteral (hydroxocobalamin or cyanocobalamin) for absorption defects or pernicious anemia, or high-dose oral (1,000--2,000 mcg daily) for dietary deficiency and food-cobalamin malabsorption. The hematologic response is monitored with reticulocyte count at 5--7 days and hemoglobin and MCV over 1--2 months; the neurological deficits improve over months but may not fully resolve. Patients with pernicious anemia require lifelong supplementation and screening for gastric neuroendocrine tumors and thyroid disease (both are associated autoimmune conditions). Serum B12 and potassium are monitored during the initial treatment phase because rapid red cell regeneration can precipitate hypokalemia.

\section*{Understanding Results}

\subsection*{Below Normal}
\begin{itemize}
  \item \textbf{Serum B12 below approximately 200 pg/mL (150 pmol/L):} Consistent with B12 deficiency, especially with compatible anemia, macrocytosis, neurological findings, or risk factors. Causes: pernicious anemia (autoimmune destruction of gastric parietal cells with antibodies to intrinsic factor and the H+/K+-ATPase), atrophic gastritis, post-gastrectomy, ileal resection or disease (Crohn disease), bacterial overgrowth, Diphyllobothrium latum (fish tapeworm) infestation, dietary deficiency (vegans, vegetarians), and increased utilization in pregnancy or malignancy.
  \item \textbf{Borderline levels (200--300 pg/mL):} Require confirmatory testing with methylmalonic acid (elevated) or homocysteine (elevated); if both are normal, deficiency is unlikely.
  \item \textbf{Clinical manifestations:} Megaloblastic anemia, hypersegmented neutrophils, glossitis, neuropsychiatric changes, subacute combined degeneration (dorsal columns and corticospinal tracts), and paresthesias. Prolonged untreated deficiency can cause irreversible neurological damage.
\end{itemize}

\subsection*{Above Normal}
\begin{itemize}
  \item \textbf{Elevated serum B12 (greater than 1,000 pg/mL):} Usually benign and reflects recent supplementation or supplementation with oral tablets, but sustained high levels can indicate increased production of binding proteins in myeloproliferative disorders (chronic myeloid leukemia, polycythemia vera, essential thrombocythemia), liver disease (hepatitis, cirrhosis), and some solid tumors.
  \item \textbf{Persistent elevation should prompt consideration of hepatic disease or a myeloproliferative neoplasm, though in most cases the finding is incidental.}
  \item \textbf{Note:} Elevated B12 does not indicate B12 excess in the sense of toxicity; B12 is water-soluble and has no known toxicity, but the elevated level is a clue to an underlying condition.
\end{itemize}

\section*{Normal Results}
\begin{itemize}
  \item \textbf{Serum B12 (adults):} 200--900 pg/mL (approximately 150--664 pmol/L), varying by laboratory and assay
  \item \textbf{Active B12 (holotranscobalamin):} greater than 35--50 pmol/L (laboratory-dependent); below 20 pmol/L strongly suggests deficiency
  \item \textbf{Methylmalonic acid (MMA):} 0.08--0.56 $\mu$mol/L in serum (varies by laboratory; elevated in B12 deficiency and renal failure)
  \item \textbf{Homocysteine:} 5--15 $\mu$mol/L (elevated in B12, folate, or B6 deficiency and in renal failure)
  \item \textbf{Daily requirement:} 2.4 mcg/day for adults (2.6 mcg in pregnancy, 2.8 mcg while lactating)
  \item \textbf{Conversion:} 1 pg/mL B12 = 0.738 pmol/L
\end{itemize}

\section*{Follow-up Tests}
\begin{itemize}
  \item \textbf{Methylmalonic acid (MMA):} The most sensitive and specific marker for B12 deficiency; elevated in more than 95\% of clinically deficient patients. Useful when serum B12 is borderline or normal but clinical suspicion is high
  \item \textbf{Homocysteine:} Elevated in B12 and folate deficiency; sensitive but less specific than MMA (also elevated in renal failure and B6 deficiency)
  \item \textbf{Folate (red blood cell and serum):} Distinguishes B12 deficiency from folate deficiency because treatment differs; both can produce megaloblastic anemia
  \item \textbf{Intrinsic factor antibodies:} Confirm the diagnosis of pernicious anemia when positive (present in 50--70\% of cases); parietal cell antibodies are more sensitive but less specific
  \item \textbf{Complete blood count with reticulocyte count and peripheral smear:} Document macrocytosis, hypersegmented neutrophils, and the response to treatment (reticulocytosis at 5--7 days)
  \item \textbf{Iron studies and ferritin:} Iron deficiency frequently coexists with B12 deficiency; replacement may be needed after B12 is initiated
  \item \textbf{Lactate dehydrogenase and indirect bilirubin:} Elevated in megaloblastic anemia from ineffective erythropoiesis and hemolysis
  \item \textbf{Upper endoscopy with gastric biopsy:} To evaluate atrophic gastritis or gastric cancer in patients with suspected pernicious anemia
  \item \textbf{Serum gastrin:} Elevated in atrophic gastritis and pernicious anemia
\end{itemize}

\section*{References}
\begin{enumerate}
  \item Stabler SP. Clinical practice: vitamin B12 deficiency. N Engl J Med. 2013;368(2):149--160.
  \item Green R, Allen LH, Bjorke-Monsen AL, et al. Vitamin B12 deficiency. Nat Rev Dis Primers. 2017;3:17040.
  \item Devalia V, Hamilton MS, Molloy AM. Guidelines for the diagnosis and treatment of cobalamin and folate disorders. Br J Haematol. 2014;166(4):496--513.
  \item Carmel R. Diagnosis and management of clinical and subclinical cobalamin deficiency: advances in 1990s. Semin Hematol. 2000;37(1):10--26.
  \item Wolffenbuttel BHR, Wouters HJCM, Heiner-Fokkema MR, van der Klauw MM. The many faces of cobalamin (vitamin B12) deficiency. Mayo Clin Proc Innov Qual Outcomes. 2019;3(2):200--214.
\end{enumerate}

\clearpage
